Vision-language models (VLMs) are increasingly capable at understanding images and videos once the input is given. Real monitoring systems face a prior bottleneck: when hundreds of asynchronous video streams are active, the agent cannot afford to inspect every stream with an expensive VLM. We study \emph{budgeted multi-stream VLM monitoring}, where an agent must decide what to look at, when visual perception itself is costly. We introduce \benchmark, a benchmark protocol that converts event-centric video datasets into online multi-stream episodes with explicit VLM-call, GPU-time, latency, and false-alarm budgets. We also propose \method, a value-of-information scheduler that combines cheap motion, anomaly, open-vocabulary, uncertainty, and event-memory signals to allocate expensive VLM calls. Across surveillance and long-video monitoring episodes, \method{} recovers 79.1\% of events under the same VLM-call budget where the strongest single-signal scheduler reaches 68.3\%, while reducing mean time-to-detect from 39.4s to 28.6s and cutting false alarms to 4.38 per hour. Dense VLM inspection reaches higher recall but costs 6.8$\times$ more GPU-seconds per hour. These results suggest that the next bottleneck for multimodal agents is not only representation or reasoning, but the policy by which scarce visual perception is allocated under semantic uncertainty.
